\documentclass[12pt,a4paper]{article}

% ------------------------------------------------
% Pacotes
% ------------------------------------------------
\usepackage[utf8]{inputenc}
\usepackage[T1]{fontenc}
\usepackage[brazil]{babel}
\usepackage{graphicx}
\usepackage[table]{xcolor}
\usepackage{array}
\usepackage{tabularx}
\usepackage{geometry}
\usepackage{fancyhdr}
\usepackage{mathptmx} % Times New Roman
\usepackage{enumitem}\usepackage{amsfonts}
\usepackage{amsmath}
\usepackage{amssymb}
\usepackage{tikz}
\usetikzlibrary{positioning, shapes.geometric}
\usepackage{listings}
\usepackage{tcolorbox}
\tcbuselibrary{breakable, skins}
 
\lstset{
    language=Java,
    basicstyle=\small\ttfamily,
    keywordstyle=\bfseries,
    showstringspaces=false,
    breaklines=true,
    frame=none,
    numbers=none,
    tabsize=2,
    commentstyle=\itshape\color{gray},
    morekeywords={this}
}

% Caixa de resolução
\tcbset{
    respostabox/.style={
        colback=green!5!white,
        colframe=green!50!black,
        fonttitle=\bfseries,
        title=Resolução,
        breakable,
        enhanced,
        before skip=6pt,
        after skip=6pt,
    }
}

\geometry{
    a4paper,
    left=2cm,
    right=2cm,
    top=2cm,
    bottom=2cm,
    headheight=15pt
}

\pagestyle{fancy}
\fancyhf{}
\renewcommand{\headrulewidth}{0pt}

\setlength{\parindent}{0pt}

% ------------------------------------------------
% Comandos auxiliares
% ------------------------------------------------
\newcolumntype{L}[1]{>{\raggedright\arraybackslash}p{#1}}
\newcolumntype{C}[1]{>{\centering\arraybackslash}p{#1}}

% ------------------------------------------------
% Documento
% ------------------------------------------------
\begin{document}

% ------------------------------------------------
% Cabeçalho
% ------------------------------------------------
\begin{tabularx}{\textwidth}{@{} L{4cm} X @{}}
    \raisebox{-0.5cm}{\includegraphics[width=4cm]{Universidade_de_Rio_Verde.jpg}} &
    \hfill
    \begin{tabular}{|>{\centering\arraybackslash}m{10cm}|>{\raggedright\arraybackslash}p{1.8cm}|}
        \hline
        \cellcolor{black}\color{white}\rule{0pt}{1.2cm}\textbf{GABARITO — AVALIAÇÃO 1}
        &
        \scriptsize Nota \tabularnewline
        \hline
    \end{tabular}
\end{tabularx}
\vspace{0.3cm}

\begin{tabularx}{\textwidth}{|X|L{2cm}|}
\hline
\textbf{Curso} & \textbf{Campus} \\
\textit{Engenharia de Software} & \textit{Rio Verde} \\
\hline
\multicolumn{2}{|l|}{ \textit{Programação Orientada a Objetos}} \\
\hline
\textbf{Professor} & \textbf{Valor} \\
\textit{Sergio Souza Novak} & 7,0 \\
\hline
\textbf{Data da Avaliação} & \textbf{Código} \\
\textit{20/03/2026} & \textit{ESW448} \\
\hline
\end{tabularx}

\vspace{0.4cm}

\begin{center}
    \rule{\textwidth}{0.4pt}
    \textbf{QUESTÕES COM RESOLUÇÕES}
\end{center}

\vspace{0.3cm}

\begin{enumerate}

% ============================================================
% QUESTÃO 1
% ============================================================
\item \textbf{(ENADE 2017)}: O encapsulamento é um mecanismo da programação orientada a objetos no qual os membros de uma classe (atributos e métodos) constituem uma caixa preta. O nível de visibilidade dos membros pode ser definido pelos modificadores de visibilidade ``privado'', ``público'' e ``protegido''. Com relação ao comportamento gerado pelos modificadores de visibilidade, assinale a opção correta.
    \begin{enumerate}[label=\Alph*)]
        \item Um atributo privado pode ser acessado pelos métodos privados da própria classe e pelos métodos protegidos das suas classes descendentes.
        \item Um atributo privado pode ser acessado pelos métodos públicos da própria classe e pelos métodos públicos das suas classes descendentes.
        \item Um membro público é visível na classe à qual ele pertence, mas não é visível nas suas classes descendentes.
        \item Um método protegido não pode acessar os atributos privados e declarados na própria classe.
        \item[\textbf{E)}] \textbf{Um membro protegido é visível na classe à qual pertence e em suas classes descendentes.} \checkmark
    \end{enumerate}

\begin{tcolorbox}[respostabox]
\textbf{Gabarito: E}

Análise das alternativas:
\begin{itemize}
    \item \textbf{A — Falsa:} Um atributo \texttt{private} \textbf{não} pode ser acessado por subclasses, nem mesmo por seus métodos \texttt{protected}. O acesso \texttt{private} é restrito à própria classe.
    \item \textbf{B — Falsa:} Atributo \texttt{private} \textbf{não} é acessível por descendentes, independentemente do modificador do método da subclasse.
    \item \textbf{C — Falsa:} Um membro \texttt{public} é visível em qualquer lugar, inclusive nas classes descendentes.
    \item \textbf{D — Falsa:} Um método \texttt{protected} \textbf{pode} acessar atributos \texttt{private} da \textit{própria} classe normalmente.
    \item \textbf{E — Verdadeira:} O modificador \texttt{protected} garante visibilidade na própria classe \textbf{e} em todas as suas subclasses (classes descendentes), mesmo em pacotes diferentes.
\end{itemize}
\end{tcolorbox}

\vspace{0.4cm}

% ============================================================
% QUESTÃO 2
% ============================================================
\item \textbf{Encapsulamento:} Implementação da classe \textbf{ContaBancaria} com o atributo \textbf{saldo} privado e os métodos \texttt{depositar}, \texttt{sacar}, \texttt{getSaldo} e construtor.

\begin{tcolorbox}[respostabox]
\begin{lstlisting}
public class ContaBancaria {
    private double saldo;

    // Construtor
    public ContaBancaria(double saldoInicial) {
        if (saldoInicial >= 0) {
            this.saldo = saldoInicial;
        } else {
            this.saldo = 0;
        }
    }

    // Deposita um valor na conta
    public void depositar(double valor) {
        if (valor > 0) {
            this.saldo += valor;
        }
    }

    // Saca um valor da conta, se houver saldo
    public void sacar(double valor) {
        if (valor > 0 && valor <= this.saldo) {
            this.saldo -= valor;
        }
    }

    // Retorna o saldo atual
    public double getSaldo() {
        return this.saldo;
    }
}
\end{lstlisting}

\textbf{Justificativa:}
\begin{itemize}
    \item O atributo \texttt{saldo} é declarado \texttt{private}, garantindo encapsulamento — nenhuma classe externa pode modificá-lo diretamente.
    \item O construtor recebe o saldo inicial e valida que não seja negativo.
    \item \texttt{depositar} e \texttt{sacar} validam os valores antes de alterar o saldo, protegendo a integridade do objeto.
    \item \texttt{getSaldo()} é o único ponto de leitura do saldo (getter), mantendo o controle sobre o acesso ao dado.
\end{itemize}
\end{tcolorbox}

\vspace{0.4cm}

% ============================================================
% QUESTÃO 3
% ============================================================
\item Na programação orientada a objetos, os modificadores de acesso controlam a visibilidade de classes, atributos e métodos. Assinale a alternativa que apresenta o modificador que permite que um membro seja acessado apenas dentro da própria classe e por suas subclasses:
    \begin{enumerate}[label=\Alph*)]
        \item \texttt{public}
        \item \texttt{private}
        \item[\textbf{C)}] \textbf{\texttt{protected}} \checkmark
        \item \texttt{package-private}
        \item \texttt{static}
    \end{enumerate}

\begin{tcolorbox}[respostabox]
\textbf{Gabarito: C — \texttt{protected}}

\begin{itemize}
    \item \texttt{public}: acessível por \textbf{qualquer} classe.
    \item \texttt{private}: acessível \textbf{apenas} dentro da própria classe (subclasses \textbf{não} têm acesso).
    \item \textbf{\texttt{protected}}: acessível dentro da própria classe \textbf{e} por suas subclasses (também acessível dentro do mesmo pacote em Java). É o gabarito correto.
    \item \texttt{package-private} (sem modificador em Java): acessível apenas dentro do mesmo pacote, mas \textbf{não} por subclasses em outros pacotes.
    \item \texttt{static}: não é um modificador de acesso; define que o membro pertence à classe, não ao objeto.
\end{itemize}
\end{tcolorbox}

\vspace{0.4cm}

% ============================================================
% QUESTÃO 4
% ============================================================
\item \textbf{Construtores:} Qual é a principal finalidade de um método construtor em uma classe?
    \begin{enumerate}[label=\Alph*)]
        \item Destruir o objeto quando ele não for mais necessário.
        \item[\textbf{B)}] \textbf{Inicializar os atributos de um objeto no momento de sua criação.} \checkmark
        \item Permitir que o objeto seja herdado por outras classes.
        \item Converter um objeto em uma string para exibição.
        \item Impedir que o objeto seja clonado.
    \end{enumerate}

\begin{tcolorbox}[respostabox]
\textbf{Gabarito: B}

O construtor é um método especial chamado automaticamente no momento em que um objeto é criado com o operador \texttt{new}. Sua principal finalidade é \textbf{inicializar os atributos do objeto}, colocando-o em um estado válido e consistente desde o início de sua vida útil.

Exemplo: \texttt{ContaBancaria c = new ContaBancaria(500.0);} — o construtor garante que o saldo começa em 500,0.

As demais alternativas descrevem:
\begin{itemize}
    \item \textbf{A:} função do \textit{Garbage Collector} (destruidor), não do construtor.
    \item \textbf{C:} função da herança (\texttt{extends}), não do construtor.
    \item \textbf{D:} função do método \texttt{toString()}.
    \item \textbf{E:} sem relação com o construtor.
\end{itemize}
\end{tcolorbox}

\vspace{0.4cm}

% ============================================================
% QUESTÃO 5
% ============================================================
\item \textbf{Indústria 4.0 — Máquina CNC:} Analise o código da classe \texttt{MaquinaCNC} e julgue as afirmações:

( ) Se a linha 14 for removida, a máquina poderá produzir peças mesmo estando desligada. \\
( ) Os atributos das linhas 03 a 08 estão corretamente encapsulados como \texttt{private}. \\
( ) O método \texttt{precisaManutencao} expõe diretamente o atributo \texttt{horasOperacao}, violando o encapsulamento.

    \begin{enumerate}[label=\Alph*)]
        \item V - V - V - F.
        \item V - F - V - V.
        \item F - V - F - V.
        \item[\textbf{D)}] \textbf{V - V - F - F.} \checkmark
        \item F - F - V - V.
    \end{enumerate}

\begin{tcolorbox}[respostabox]
\textbf{Gabarito: D — V - V - F}

Análise de cada afirmação:

\textbf{1ª afirmação — VERDADEIRA:} A linha 14 é \texttt{if (!this.ligada)\{ return false; \}}, que impede produção quando a máquina está desligada. Se for removida, o método \texttt{produzir} executará mesmo com \texttt{ligada = false}.

\textbf{2ª afirmação — VERDADEIRA:} Todos os atributos das linhas 03–08 (\texttt{modelo}, \texttt{horasOperacao}, \texttt{temperaturaAtual}, \texttt{temperaturaMaxima}, \texttt{ligada}, \texttt{pecasProduzidas}) são declarados \texttt{private}, o que os encapsula corretamente.

\textbf{3ª afirmação — FALSA:} O método \texttt{precisaManutencao()} é \texttt{public} e chama internamente o método \texttt{private} \texttt{verificarManutencao()}, que por sua vez acessa \texttt{horasOperacao}. O atributo \textbf{não é exposto diretamente}; o encapsulamento é preservado — apenas um resultado booleano é devolvido.

\textit{Observação:} A questão lista três afirmações, portanto a sequência tem três termos: \textbf{V - V - F}.
\end{tcolorbox}

\vspace{0.4cm}

% ============================================================
% QUESTÃO 6
% ============================================================
\item \textbf{Siderurgia — Alto Forno:} Julgue as afirmações sobre a classe \texttt{AltoForno}:

( ) O método \texttt{verificarFaixaTermica} é \texttt{private}, podendo ser invocado apenas por outras classes. \\
( ) Na linha 15, \texttt{controleSeguranca} é chamado para garantir que a temperatura não ultrapasse o limite. \\
( ) A produção de metal depende exclusivamente do forno estar \texttt{ativo}. \\
( ) O atributo \texttt{temperaturaInterna} poderia ser \texttt{public} sem afetar a segurança.

    \begin{enumerate}[label=\Alph*)]
        \item V - V - F - F.
        \item V - F - V - V.
        \item[\textbf{C)}] \textbf{F - V - F - F.} \checkmark
        \item V - V - V - F.
        \item F - F - F - V.
    \end{enumerate}

\begin{tcolorbox}[respostabox]
\textbf{Gabarito: C — F - V - F - F}

Análise de cada afirmação:

\textbf{1ª afirmação — FALSA:} \texttt{verificarFaixaTermica} é \texttt{private}, o que significa que só pode ser invocado por métodos da \textbf{própria classe} \texttt{AltoForno}. A afirmação diz ``por outras classes'', o que é o oposto — portanto FALSA.

\textbf{2ª afirmação — VERDADEIRA:} Dentro do método \texttt{injetarCombustivel} (linha 19), após incrementar a temperatura, \texttt{controleSeguranca()} é chamado (linha 21). Esse método desliga o forno se a temperatura ultrapassar o máximo, demonstrando controle de estado interno.

\textbf{3ª afirmação — FALSA:} O método \texttt{produzirMetal} (linha 25–32) verifica \textbf{duas} condições: \texttt{ativo} \textbf{E} \texttt{verificarFaixaTermica()}. A temperatura deve estar dentro da faixa operacional — portanto não depende \textit{exclusivamente} de \texttt{ativo}.

\textbf{4ª afirmação — FALSA:} Tornar \texttt{temperaturaInterna} \texttt{public} permitiria que qualquer classe externa alterasse a temperatura arbitrariamente, driblando o \texttt{controleSeguranca}. Isso comprometeria a integridade do objeto e do equipamento.
\end{tcolorbox}

\vspace{0.4cm}

% ============================================================
% QUESTÃO 7
% ============================================================
\item \textbf{(VUNESP 2021)} Em programação orientada a objetos, métodos de acesso do tipo \textit{setter} têm a finalidade primária de:
    \begin{enumerate}[label=\Alph*)]
        \item obter o valor de um atributo.
        \item[\textbf{B)}] \textbf{modificar o valor de um atributo.} \checkmark
        \item inicializar os valores de atributos de um objeto assim que ele é instanciado.
        \item contar quantas vezes o valor de um atributo foi obtido.
        \item salvar o valor de um atributo de um objeto em um banco de dados.
    \end{enumerate}

\begin{tcolorbox}[respostabox]
\textbf{Gabarito: B}

\begin{itemize}
    \item \textbf{Getter} (\texttt{get...}): obtém/retorna o valor de um atributo privado.
    \item \textbf{Setter} (\texttt{set...}): \textbf{modifica} o valor de um atributo privado, podendo incluir validações antes da alteração.
    \item \textbf{A} descreve o \textit{getter}.
    \item \textbf{C} descreve o \textit{construtor}.
    \item \textbf{D} e \textbf{E} não correspondem a nenhum padrão estabelecido de POO.
\end{itemize}
\end{tcolorbox}

\vspace{0.4cm}

% ============================================================
% QUESTÃO 8
% ============================================================
\item \textbf{(FUNDATEC 2025 — GHC-RS)} Considerando o código abaixo, quem pode acessar o atributo \texttt{v} diretamente:

\begin{lstlisting}
package conceicao;
public class Conceicao {
    protected String v;
}
\end{lstlisting}

    \begin{enumerate}[label=\Alph*)]
        \item Apenas as classes que estejam no mesmo pacote.
        \item Apenas os métodos da própria classe.
        \item[\textbf{C)}] \textbf{As classes filhas de qualquer pacote.} \checkmark
        \item As classes filhas apenas do mesmo pacote.
        \item Por qualquer classe de qualquer pacote.
    \end{enumerate}

\begin{tcolorbox}[respostabox]
\textbf{Gabarito: C}

Em Java, o modificador \texttt{protected} concede acesso:
\begin{enumerate}
    \item A todos os membros do \textbf{mesmo pacote} (\textit{package access}).
    \item A todas as \textbf{subclasses} (\texttt{extends}), \textbf{independentemente do pacote}.
\end{enumerate}

Portanto, uma classe filha em \textbf{qualquer pacote} pode acessar \texttt{v} diretamente por herança.

\textbf{Atenção:} A alternativa A está incompleta (omite as subclasses de outros pacotes). A alternativa E descreve \texttt{public}. A alternativa B descreve \texttt{private}.
\end{tcolorbox}

\vspace{0.4cm}

% ============================================================
% QUESTÃO 9 (última questão discursiva)
% ============================================================
\item Considere a modelagem da classe \texttt{Carro} com atributo privado \texttt{velocidadeAtual}. Explique a importância de \textit{getters} e \textit{setters} e dê um exemplo de regra de negócio para \texttt{setVelocidade()}.

\begin{tcolorbox}[respostabox]
\textbf{Importância dos métodos \textit{getter} e \textit{setter}:}

Quando um atributo é declarado \texttt{private}, nenhuma classe externa pode lê-lo ou alterá-lo diretamente. Os métodos \textit{getter} e \textit{setter} são a \textbf{interface controlada} para acessar e modificar esses valores. Isso traz as seguintes vantagens:

\begin{enumerate}
    \item \textbf{Controle de acesso:} somente o que for exposto via métodos públicos é visível ao mundo externo. Atributos sensíveis ficam protegidos.
    \item \textbf{Validação de dados:} o \textit{setter} pode verificar se o valor recebido é válido antes de aplicá-lo, evitando estados inconsistentes no objeto.
    \item \textbf{Manutenibilidade:} a lógica interna da classe pode mudar sem afetar o código que a utiliza, desde que a assinatura dos métodos permaneça a mesma.
    \item \textbf{Rastreabilidade:} é possível adicionar log, notificações ou outros comportamentos quando um atributo é lido ou modificado, sem alterar código externo.
\end{enumerate}

\textbf{Exemplo de regra de negócio no \texttt{setVelocidade()}:}

O método \texttt{setVelocidade()} poderia implementar a seguinte regra: \textit{``A velocidade do carro não pode ser negativa e não pode ultrapassar 200 km/h (velocidade máxima permitida pelo veículo). Se o valor fornecido for inválido, a velocidade permanece inalterada e uma mensagem de aviso é registrada.''}

Dessa forma, qualquer parte do sistema que tente definir uma velocidade de $-10$ ou de $300\,\text{km/h}$ será bloqueada pela própria classe, garantindo que o objeto \texttt{Carro} nunca fique em um estado inconsistente com a realidade.
\end{tcolorbox}

\end{enumerate}

\vspace{0.8cm}

\end{document}
